\chapter{An Introduction to Exact Categories} \label{chapter:exact}

\section{Motivation}
When studying abelian categories, there are many convenient properties related to monics and epics. Let us recall a few:

\begin{prop} \label{prop:popb_abelian}
    In an abelian category, the following properties hold:
    \begin{itemize}
        \item the pushout of a monic is a monic;
        \item the pullback of a monic is a monic;
        \item the pushout of an epic is an epic;
        \item the pullback of an epic is an epic.
    \end{itemize}
\end{prop}

Another nice property is the following:

\begin{prop} \label{prop:kercoker}
    In an abelian category,
    \begin{itemize}
        \item every monic is a kernel of its cokernel;
        \item every epic is a cokernel of its kernel.
    \end{itemize}
\end{prop}

This allows us to conclude many properties about kernels and cokernels:

\begin{prop} \label{prop:ab_is_exact}
    In an abelian category, the class of kernels and the class of cokernels are closed under isomorphism. Furthermore, the following hold:
    \begin{itemize}[align=right]
        \item[\textnormal{[E0]}] For any object $C \in \cat{C}$, the identity morphism $1_C$ is a kernel;
        \item[\textnormal{[E0$^{\op}$]}] For any object $C \in \cat{C}$, the identity morphism $1_C$ is a cokernel;
        \item[\emph{[E1]}] The class of kernels is closed under composition;
        \item[\emph{[E1$^{\op}$]}] The class of cokernels is closed under composition;  
        \item[\emph{[E2]}] The pushout of a kernel along any morphism is a kernel;
        \item[\emph{[E2$^{\op}$]}] The pullback of a cokernel along any morphism is a cokernel.
    \end{itemize}
\end{prop}

\begin{preuve}
    It is clear that kernels and cokernels are closed under isomorphism.

    E0 and E0$^{\op}$ follow directly from the fact that the identity morphism is both the kernel and cokernel of the zero map.

    We can deduce E1 and E1$^{\op}$ by recalling that composing monics (respectively epics) gives us monics and that monics are kernels of their cokernel (respectively epics are cokernels of their kernel).

    As for E2, recall that kernels are always monic. Thus, by \ref{prop:popb_abelian}, the pushout of a kernel is monic as well. But then, by \ref{prop:kercoker}, it is a kernel of its cokernel, thus it is a kernel.
    The proof for E2$^{\op}$ is similar.
\end{preuve}

\section{Exact Categories}
Based on the observations from the last section, one can define a more general notion that encapsulates the idea of short exact sequences, namely exact categories. To do so, let us first define kernel-cokernel pairs.

\begin{defn}
    Let $\cat{A}$ be an additive category. A \emph{kernel-cokernel pair} $(i,p)$ in $\cat{A}$ is a pair of composable morphisms $A \xra{i} B \xra{p} C$ such that $i$ is a kernel of $p$ and $p$ is a cokernel of $i$.
\end{defn}

In the context of abelian categories, we notice that kernel-cokernel pairs indeed give rise to short exact sequences in the usual sense. In a general additive category, however, we cannot say the same since we have no guarantee that kernels and cokernels even exist for arbitrary morphisms, and thus we are unsure that $\Img{i}$ exists. But we will use our intuition from abelian categories to expand the notion of short exact sequences to additive categories.

\begin{defn}
    Let $\cat{A}$ be an additive category and $\cat{E}$ a class of kernel-cokernel pairs on $\cat{A}$. We say that $A \xra{i} B$ is an \emph{admissible monic} if there is a morphism $B \xra{p} C$ in $\cat{A}$ such that $(i,p) \in \cat{E}$. Dually, $B \xra{q} C$ is an \emph{admissible epic} if there exists a morphism $A \xra{j} B$ such that $(j,q) \in \cat{E}$.

    A class $\cat{E}$ of kernel-cokernel pairs is called an \emph{exact structure} on $\cat{A}$ if it is closed under isomorphisms and satisfies the following axioms:
    \begin{itemize}[align=right]
        \item[{[E0]}] For any object $A \in \cat{A}$, the identity morphism $1_A$ is an admissible monic;
        \item[{[E0$^{\op}$]}] For any object $A \in \cat{A}$, the identity morphism $1_A$ is an admissible epic;
        \item[{[E1]}] The class of admissible monics is closed under composition;
        \item[{[E1$^{\op}$]}] The class of admissible epics is closed under composition;
        \item[{[E2]}] The pushout of an admissible monic along any morphism exists and is an admissible monic;
        \item[{[E2$^{\op}$]}] The pullback of an admissible epic along any morphism exists and is an admissible epic.
    \end{itemize}
    An \emph{exact category} is a pair $(\cat{A},\cat{E})$ where $\cat{A}$ is an additive category and $\cat{E}$ is an exact structure on $\cat{A}$. In that context, elements of $\cat{E}$ are called \emph{short exact sequences}.
\end{defn}



